\section*{3 Binary Tree Codes:Four Colors Suffice}
In this section we prove Theorem2. We start by setting some notation. Let $T$ be the infinite complete rooted binary tree. We identify length-$n$ paths in $T$ that starts at the root with length-$n$ binary strings in the natural way. Namely, we identify left son and right son with 0 and 1, respectively. Given a node $v$ at depth $n\ge1$ we define $p_v\in\{0,1\}^n$ to be the string that encodes the (unique) path from the root to $v$.
An edge-coloring of $T$ by a color set $\Sigma$ is given by a function, which for ease of readability, we slightly abuse notation and also denote by $T:\{0,1\}^{\mathbb N_1}\to\Sigma^{\mathbb N_1}$, where the color of an edge $e=u\to v$ is $T(p_v)_{\operatorname{depth}(v)}$. Note that $T$ is an online function, namely, for every $x\in\{0,1\}^{\mathbb N_1}$ and $i\in\mathbb N_1$, the value $T(x)_i$ is determined by $x_1,\ldots,x_i$.
\paragraph{The (probabilistic) construction}
Let $\{R_i\}_{i\in\mathbb N_1}$ be a sequence of independent random variables, each is uniformly distributed over $\mathbb F_4$--the field of 4 elements. Let $\mathbb F_2$ be the (unique) subfield of $\mathbb F_4$ of size 2. Define the (random) coloring function $T:\mathbb F_2^{\mathbb N_1}\to\mathbb F_4^{\mathbb N_1}$ (where we identify $\mathbb F_2$ and $\{0,1\}$ in the natural way) as follows: for every $t\in\mathbb N_1$
\begin{equation}
T(x)_t=\sum_{i=1}^tR_{t+1-i}x_i.\tag{1}
\end{equation}
\begin{authorclaim}{Definition9}
Let $v$ be a depth-$n$ vertex in $T$. Let $\ell\ge1$ and $x,y\in\mathbb F_2^{\ell-1}$. For $k=1,\ldots,\ell$ we define the random variable
\[
a_v(x,y,k)=T(p_v\circ1\circ y)_{n+k}-T(p_v\circ0\circ x)_{n+k}.
\]
Note that $a_v(x,y,k)$ is a (random) element in $\mathbb F_4$. We define the integral random variable $h_v(x,y)=\sum_{k=1}^\ell I_k$, where $I_k$ is the indicator random variable that equals 1 when $a_v(x,y,k)\ne0$. Note that $h_v(x,y)\in\{0,1,\ldots,\ell\}$ is the Hamming distance between $T(p_v\circ0\circ x)_{[n+1,n+\ell]}$ and $T(p_v\circ1\circ y)_{[n+1,n+\ell]}$.
\end{authorclaim}
\begin{authorclaim}{Claim10}
Let $v$ be a vertex in $T$. Let $\ell\ge1$ and $x,y\in\mathbb F_2^{\ell-1}$. Then, for every $k\in\{1,\ldots,\ell\}$ it holds that $a_v(x,y,k)=R_k+\sum_{i=1}^{k-1}R_{k-i}(y-x)_i$.
\end{authorclaim}
\begin{proof}
Denote the depth of $v$ by $n$. Fix $k\in\{1,\ldots,\ell\}$. By Equation(1),
\begin{align*}
T(p_v\circ0\circ x)_{n+k}
&=\sum_{i=1}^{n+k}R_{n+k+1-i}(p_v\circ0\circ x)_i\\
&=\sum_{i=1}^nR_{n+k+1-i}(p_v)_i+\sum_{i=1}^kR_{k+1-i}(0\circ x)_i.
\end{align*}
Similarly
\[
T(p_v\circ1\circ y)_{n+k}
=\sum_{i=1}^nR_{n+k+1-i}(p_v)_i+\sum_{i=1}^kR_{k+1-i}(1\circ y)_i.
\]
Thus,
\begin{align*}
a_v(x,y,k)
&=\sum_{i=1}^kR_{k+1-i}(1\circ y)_i-\sum_{i=1}^kR_{k+1-i}(0\circ x)_i\\
&=R_k+\sum_{i=1}^{k-1}R_{k-i}(y-x)_i.
\end{align*}
\end{proof}
\begin{authorclaim}{Claim11}
Let $v$ be a vertex in $T$. Let $\ell\ge1$ and $x,y\in\mathbb F_2^{\ell-1}$. Then, the random variables $a_v(x,y,1),\ldots,a_v(x,y,\ell)$ are independent and each is uniformly distributed over $\mathbb F_4$.
\end{authorclaim}
\begin{proof}
By Claim10, $a_v(x,y,k)=R_k+L_k$ where $L_k$ is some $\mathbb F_4$-linear combination of $R_1,\ldots,R_{k-1}$. Therefore, $a_v(x,y,k)$ is independent of the joint distribution of $a_v(x,y,1),\ldots,a_v(x,y,k-1)$. As this holds for every $k$ we have that $a_v(x,y,1),\ldots,a_v(x,y,\ell)$ are independent. To conclude the proof, note that for every fixing of $R_1,\ldots,R_{k-1}$, $a_v(x,y,k)=R_k+\ell_k$ for some fixed $\ell_k\in\mathbb F_4$ and so $a_v(x,y,k)$ is uniform over $\mathbb F_4$.
\end{proof}
\begin{authorclaim}{Claim12}
For every two vertices $u,v$ in $T$ and every $x,y\in\mathbb F_2^{\ell-1}$, $h_v(x,y)=h_u(x,y)$, and $h_v(x,y)=h_v(0^{\ell-1},y-x)$.
\end{authorclaim}
\begin{proof}
The first equality follows immediately by Claim10 as, for every $k\in\{1,\ldots,\ell\}$, the expression obtained for $a_v(x,y,k)$ is independent of the choice of $v$. As for the second asserted equality, again by Claim10,
\begin{align*}
a_v(x,y,k)&=R_k+\sum_{i=1}^{k-1}R_{k-i}(y-x)_i\\
&=R_k+\sum_{i=1}^{k-1}R_{k-i}((y-x)-0)_i
=a_v(0^{\ell-1},y-x,k),
\end{align*}
where observe that for the last equality we are using the fact that $\mathbb F_2$ is a subfield of $\mathbb F_4$ and so $y-x\in\mathbb F_2^{\ell-1}$. Indeed, recall that $a_v$'s second argument is a binary string and so the equality above would have been meaningless otherwise. The above equation implies $h_v(x,y)=h_v(0^{\ell-1},y-x)$, proving the claim.
\end{proof}
Given Claim12 we can simplify our notation as follows. Let $r$ denote the root of $T$. For $x\in\{0,1\}^{\ell-1}$ and $k\in\{1,\ldots,\ell\}$ we define the random variables $a(x,k)=a_r(0^\ell,1\circ x,k)$, $h(x)=h_r(0^{\ell-1},x)$. Note that $h(x)=\sum_{k=1}^\ell a(x,k)$.
\begin{authorclaim}{Theorem13}
There exists a fixing of the sequence $\{R_i\}_i$ such that the function $T$ is a tree code with distance 0.05.
\end{authorclaim}
\begin{proof}
First note that for every fixing of the sequence $\{R_i\}_i$, $T$ is an online function. Observe that, for a fixing of $\{R_i\}_i$, $T$ is a tree code with distance $\delta$ if and only if for every $\ell\ge1$ and $x\in\{0,1\}^{\ell-1}$ it holds that $h(x)\ge\delta\ell$. Indeed, recall that by definition, $T$ is a tree code with distance $\delta$ if and only if for every vertex $v$ in $T$, $\ell\ge1$, and for every $x,y\in\{0,1\}^{\ell-1}$ it holds that $h_v(x,y)\ge\delta\ell$. However, by Claim12, $h_v(x,y)=h(y-x)$.
For $x\in\{0,1\}^{\ell-1}$ denote by $E(x)$ the event $h(x)<\delta\ell$. By the above discussion, it suffices to prove, for $\delta=0.05$, that $\Pr[\bigcup_{x\in\{0,1\}^{\mathbb N}}E(x)]<1$. To this end, by the union bound, it suffices to prove that $\sum_{x\in\{0,1\}^{\mathbb N}}\Pr[E(x)]<1$.
Consider any $x\in\{0,1\}^{\ell-1}$ with $\ell\ge1$. Note that the event $E(x)$ holds if and only if there exists a set $\mathcal T\subseteq\{1,\ldots,\ell\}$ of size $|\mathcal T|\ge\lceil(1-\delta)\ell\rceil$ such that for every $k\in\mathcal T$, $a(x,k)=0$. By taking the union bound over all such sets $\mathcal T$, and using that $a(x,1),\ldots,a(x,\ell)$ are independent and each is uniformly distributed over $\mathbb F_4$ (Claim11), we get
\begin{equation}
\Pr[E(x)]\le\binom\ell{\lceil(1-\delta)\ell\rceil}4^{-\lceil(1-\delta)\ell\rceil}.\tag{2}
\end{equation}
By Lemma7, we have that
\[
\frac1\ell\log_2\binom\ell{\lceil(1-\delta)\ell\rceil}
\le H\left(\frac{\lceil(1-\delta)\ell\rceil}{\ell}\right).
\]
As $\delta<1/2$ and since the entropy function $H$ decreases in $[1/2,1]$ we have that
$H(\lceil(1-\delta)\ell\rceil/\ell)\le H(1-\delta)=H(\delta)$.
Substitute to Equation(2), we get that $\Pr[E(x)]\le2^{(H(\delta)-2(1-\delta))\ell}$. Thus,
\begin{align*}
\sum_{x\in\{0,1\}^{\mathbb N}}\Pr[E(x)]
&\le\sum_{\ell=1}^\infty2^{\ell-1}\cdot2^{(H(\delta)-2(1-\delta))\ell}\\
&=\frac12\sum_{\ell=1}^\infty2^{(H(\delta)+2\delta-1)\ell}.
\end{align*}
One can verify that for $\delta=0.05$ the above geometric sum is strictly smaller than1, and the theorem follows.
\end{proof}
\subsection*{3.1 Improving the distance}
We now show a method for improving the distance. We illustrate it to obtain a bound of 0.136 on the distance, which proves Theorem2, though we believe that the method can be used to push the bound further. It is fairly easy to show that the distance of a 4-color tree code cannot be larger than 1/2.
\begin{authorclaim}{Theorem14}
There exists a fixing of the sequence $\{R_i\}_i$ such that the function $T$ is a tree code with distance 0.136.
\end{authorclaim}
\begin{proof}
For the proof it will be convenient to consider a specific representation of $\mathbb F_4$. We make use of the standard construction of $\mathbb F_4$ as a quotient of the polynomial ring over $\mathbb F_2$ with respect to an ideal generated by a degree 2 irreducible element as follows. Note that $t^2+t+1\in\mathbb F_2[t]$ is irreducible, and so $K=\mathbb F_2[t]/\langle t^2+t+1\rangle$ is a field of 4 elements which we will take as the construction for $\mathbb F_4$. Let $\alpha$ be the class of $t$ in $K$. In this representation, the field $\mathbb F_4$ consists of the elements $0,1,\alpha,\alpha+1$ where $\alpha^2+\alpha+1=0$.
Consider the sequence $\{R_i\}_{i\in\mathbb N}$ as in the beginning of the section but with the fixings $R_1=1$ and $R_2=\alpha$. Observe that for every $x\in\mathbb F_2^{\ell-1}$ with $\ell\ge2$ it holds that $a(x,1)=1$ and $a(x,2)=\alpha+x_1$. In particular, $a(x,1),a(x,2)$ are both non-zeros and so $h(x)\ge2$. Let $\ell_0\ge2$ be an integer parameter to be chosen later on. By the above, we have that for every $x\in\mathbb F_2^{\ell-1}$ with $\ell\le\ell_0$ it holds that
\begin{equation}
\frac{h(x)}{\ell}\ge\frac2{\ell_0}.\tag{3}
\end{equation}
For $x\in\{0,1\}^{\ell-1}$ denote by $E_{1,\alpha}(x)$ the event $h(x)<\delta\ell$ with the $\{R_i\}_{i\in\mathbb N}$ as defined above, namely, $R_1=1$, $R_2=\alpha$ and the rest of the random variables $\{R_i\mid i\ge3\}$ are independent and uniformly distributed over $\mathbb F_4$. Once we establish a bound of
\begin{equation}
\Pr[\bigcup_{|x|\ge\ell_0}E_{1,\alpha}(x)]<1\tag{4}
\end{equation}
for some choice of $\delta$ then, combined with Equation(3), we will establish the existence of a tree code with distance at least $\min(2/\ell_0,\delta)$.
Consider any $x\in\{0,1\}^{\ell-1}$ with $\ell\ge\ell_0+1$. The event $E_{1,\alpha}(x)$ holds if and only if there exists a set $\mathcal T\subseteq\{3,\ldots,\ell\}$ of size $|\mathcal T|\ge\lceil(1-\delta)\ell\rceil$ such that for every $k\in\mathcal T$, $a(x,k)=0$. By taking the union bound over all such sets $\mathcal T$, and using that $a(x,3),\ldots,a(x,\ell)$ are independent and each is uniformly distributed over $\mathbb F_4$, we get that
\begin{align*}
\Pr[E_{1,\alpha}(x)]
&\le\binom{\ell-2}{\lceil(1-\delta)\ell\rceil}4^{-\lceil(1-\delta)\ell\rceil}\\
&\le\binom\ell{\lceil(1-\delta)\ell\rceil}4^{-\lceil(1-\delta)\ell\rceil}.
\end{align*}
By Lemma7, we have that
\[
\frac1\ell\log_2\binom\ell{\lceil(1-\delta)\ell\rceil}
\le H\left(\frac{\lceil(1-\delta)\ell\rceil}{\ell}\right).
\]
As we will choose $\delta<1/2$ and the entropy function $H$ decreases in $[1/2,1]$ we have that $H(\lceil(1-\delta)\ell\rceil/\ell)\le H(1-\delta)=H(\delta)$.
Thus, $\Pr[E_{1,\alpha}(x)]\le2^{(H(\delta)-2(1-\delta))\ell}$. By substituting the above equation to Equation(4), we get that
\[
\sum_{|x|\ge\ell_0}\Pr[E_{1,\alpha}(x)]
\le\sum_{\ell=\ell_0+1}^\infty2^{\ell-1}\cdot2^{(H(\delta)-2(1-\delta))\ell}.
\]
Write $\beta=2^{H(\delta)+2\delta-1}$. Then, the above is bounded by $\frac12\sum_{\ell=\ell_0+1}^\infty\beta^\ell=\beta^{\ell_0+1}/(2(1-\beta))$.
Consider the real polynomial $f_{\ell_0}(x)=x^{\ell_0+1}-2(1-x)$. We have that $f'_{\ell_0}(x)=(\ell_0+1)x^{\ell_0}+2$.
Since $\ell_0\ge2$, $f'_{\ell_0}(x)>0$ for all $x\ge0$. Further, $f_{\ell_0}(0)=-2$ and $f_{\ell_0}(1)=1$. Thus, $f_{\ell_0}(x)$ has a single root $\beta_{\ell_0}\in[0,1]$ (in fact, $\beta_{\ell_0}$ is monotone-increasing as a function of $\ell_0$, and $\beta_{\ell_0}\to1$ as $\ell_0\to\infty$). For a fixed choice of $\ell_0$, by choosing $\beta<\beta_{\ell_0}$ and solving for $\delta$ (recall $\beta=2^{H(\delta)+2\delta-1}$) to obtain $\delta_{\ell_0}$, we get that there exists a fixing of $\{R_i\mid i\ge3\}$ such that the obtained tree code has distance at least $\min(\delta_{\ell_0},2/\ell_0)$. Thus, the obtained bound is $\max_{\ell_0\ge2}\min(\delta_{\ell_0},2/\ell_0)$.
One can verify that $\ell_0=14$ maximizes the above equation to get distance larger than 0.136.
\end{proof}
